\section{Trees}

\entry{Euler tour}{$O(n)$ \quad mem $n$}{%
The subtree of \code{u} is the range \code{[tin[u], tout[u])} in any array indexed by
\code{tin}, and \code{u} is an ancestor of \code{v} iff
\code{tin[u] <= tin[v] < tout[u]}. Turns ``update / query a whole subtree'' into a
range operation, so pair it with \code{BIT} / \code{Seg} / \code{LazySeg}. Edge values
go on the \emph{lower} endpoint of the edge.\\[0.3ex]
Iterative, so depth never matters.}

\begin{lstlisting}
void euler(const vector<vector<int>>& g, int root,
           vector<int>& tin, vector<int>& tout) {
    int n = (int)g.size(), t = 0;
    tin.assign(n, 0); tout.assign(n, 0);
    vector<int> par(n, -1), it(n, 0), st{root};
    tin[root] = t++;
    while (!st.empty()) {
        int u = st.back();
        if (it[u] < (int)g[u].size()) {
            int v = g[u][it[u]++];
            if (v == par[u]) continue;
            par[v] = u; tin[v] = t++; st.push_back(v);
        } else { tout[u] = t; st.pop_back(); }
    }
}
\end{lstlisting}

\entry{LCA (binary lifting)}{build $O(n \log n)$, query $O(\log n)$ \quad mem $n \log n$}{%
\code{up[k][v]} is the $2^k$-th ancestor of \code{v}; the root is its own ancestor so
jumps past it clamp there. Lift the deeper vertex to the other's depth, then jump both
up together while the ancestors differ. Built by BFS, so no recursion depth
trouble.\\[0.3ex]
Max / min / sum edge on a path: add a second table \code{mx} with \code{mx[0][v]} =
weight of the edge to the parent, \code{mx[0][root] = ID},
\code{mx[k][v] = op(mx[k-1][v], mx[k-1][up[k-1][v]])}, and merge \code{mx[i][v]} into
the answer at every jump taken in \code{kth} and \code{lca}.}

\begin{lstlisting}
struct LCA {
    int LG; vector<vector<int>> up; vector<int> dep;
    LCA(const vector<vector<int>>& g, int root = 0) {
        int n = g.size(); LG = 1 + __lg(max(n, 2));
        up.assign(LG, vector<int>(n, root)); dep.assign(n, 0);
        vector<int> order{root};
        vector<bool> seen(n); seen[root] = true;
        for (int i = 0; i < (int)order.size(); i++) {
            int u = order[i];
            for (int v : g[u]) if (!seen[v]) {
                seen[v] = true; up[0][v] = u;
                dep[v] = dep[u] + 1; order.push_back(v);
            }
        }
        for (int k = 1; k < LG; k++)
            for (int v = 0; v < n; v++)
                up[k][v] = up[k-1][up[k-1][v]];
    }
    int kth(int v, int k) {        // root if k is too big
        k = min(k, dep[v]);
        for (int i = 0; k; k >>= 1, i++)
            if (k & 1) v = up[i][v];
        return v;
    }
    int lca(int a, int b) {
        if (dep[a] < dep[b]) swap(a, b);
        a = kth(a, dep[a] - dep[b]);
        if (a == b) return a;
        down(k, LG, 0) if (up[k][a] != up[k][b])
            a = up[k][a], b = up[k][b];
        return up[0][a];
    }
    int dist(int a, int b) {
        return dep[a] + dep[b] - 2 * dep[lca(a, b)];
    }
};
\end{lstlisting}

\entry{Heavy-light decomposition}{build $O(n)$, path / subtree $O(\log^2 n)$ \quad mem $n$}{%
Every vertex points at its heaviest child, so a root-to-leaf walk changes chain at most
$\log n$ times. Positions are handed out so that each chain and each subtree is a
contiguous range; feed those to \code{Seg} / \code{LazySeg} / \code{BIT}, which is the
second log. \code{lca} comes for free, no extra table.\\[0.3ex]
\code{path} calls \code{f} on $O(\log n)$ ranges in \emph{arbitrary} order, so
\code{op} must commute. For a non-commutative op keep two trees, one on \code{pos} and
one on reversed \code{pos}, and collect the two sides of the LCA separately.
\code{edges = true} stores each edge at its lower endpoint and skips the LCA itself;
then \code{sub} should start at \code{pos[u] + 1} too. Recursion depth is $O(n)$ ---
on a path of $2\cdot10^5$ vertices raise the stack or rewrite the two \code{dfs}
calls iteratively.\\[0.3ex]
\code{HLD h(g); Seg s(n); s.set(h.pos[u], val[u]);}\\
\code{h.path(a, b, [\&](int l, int r)\{ ... s.query(l, r) ... \});}}

\begin{lstlisting}
struct HLD {
    int n, t = 0;
    vector<vector<int>> g;      // heavy child kept first
    vector<int> par, dep, sz, head, pos;
    HLD(vector<vector<int>> adj, int root = 0)
        : n(adj.size()), g(move(adj)), par(n, -1), dep(n),
          sz(n), head(n, root), pos(n) {
        dfs(root); build(root, root);
    }
    void dfs(int u) {
        if (par[u] != -1) g[u].erase(find(all(g[u]), par[u]));
        sz[u] = 1;
        for (int& v : g[u]) {
            par[v] = u; dep[v] = dep[u] + 1;
            dfs(v); sz[u] += sz[v];
            if (sz[v] > sz[g[u][0]]) swap(v, g[u][0]);
        }
    }
    void build(int u, int h) {      // heavy child keeps head
        head[u] = h; pos[u] = t++;
        for (int i = 0; i < (int)g[u].size(); i++)
            build(g[u][i], i ? g[u][i] : h);
    }
    int lca(int a, int b) {
        for (; head[a] != head[b]; b = par[head[b]])
            if (dep[head[a]] > dep[head[b]]) swap(a, b);
        return dep[a] < dep[b] ? a : b;
    }
    pair<int,int> sub(int u) {              // subtree range
        return {pos[u], pos[u] + sz[u]};
    }
    template<class F>
    void path(int a, int b, F f, bool edges = false) {
        for (; head[a] != head[b]; b = par[head[b]]) {
            if (dep[head[a]] > dep[head[b]]) swap(a, b);
            f(pos[head[b]], pos[b] + 1);
        }
        if (dep[a] > dep[b]) swap(a, b);
        f(pos[a] + edges, pos[b] + 1);
    }
};
\end{lstlisting}

\entry{Tree diameter}{$O(n)$ \quad mem $n$}{%
Two BFS runs: the farthest vertex from anywhere is an endpoint of some diameter, and
the farthest vertex from \emph{that} is the other endpoint. Unweighted only --- with
weights use the same shape with Dijkstra, or a single DFS keeping the two deepest
children per vertex. Iterative.}

\begin{lstlisting}
// returns {length in edges, endpoint a, endpoint b}
array<int,3> diameter(const vector<vector<int>>& g) {
    int n = (int)g.size();
    auto bfs = [&](int s) {
        vector<int> d(n, -1); d[s] = 0;
        vector<int> q{s};
        for (int i = 0; i < (int)q.size(); i++)
            for (int v : g[q[i]]) if (d[v] < 0) {
                d[v] = d[q[i]] + 1; q.push_back(v);
            }
        int b = s;
        for (int i = 0; i < n; i++)
            if (d[i] > d[b]) b = i;
        return pair<int,int>{b, d[b]};
    };
    int a = bfs(0).first;
    auto [b, len] = bfs(a);
    return {len, a, b};
}
\end{lstlisting}

\entry{Centroid}{$O(n)$ \quad mem $n$}{%
The centroid minimises the largest component left after deleting it; that component is
always $\le n/2$. Every tree has one or two, and two only when they are
adjacent.\\[0.3ex]
Recurse on the pieces (removing the centroid each time) and you get centroid
decomposition: $O(\log n)$ levels, every path in the tree passes through the centroid
of exactly one level, which is how ``count paths with property X'' problems are
solved in $O(n \log n)$.}

\begin{lstlisting}
vector<int> centroids(const vector<vector<int>>& g) {
    int n = (int)g.size();
    vector<int> sub(n, 1), par(n, -1), it(n, 0), st{0}, ord;
    while (!st.empty()) {              // iterative dfs order
        int u = st.back();
        if (it[u] < (int)g[u].size()) {
            int v = g[u][it[u]++];
            if (v == par[u]) continue;
            par[v] = u; st.push_back(v);
        } else { ord.push_back(u); st.pop_back(); }
    }
    for (int u : ord)                        // children first
        if (par[u] >= 0) sub[par[u]] += sub[u];
    vector<int> res;
    for (int u = 0; u < n; u++) {
        int mx = n - sub[u];
        for (int v : g[u]) if (v != par[u])
            mx = max(mx, sub[v]);
        if (mx <= n / 2) res.push_back(u);
    }
    return res;
}
\end{lstlisting}

\entry{Rerooting DP}{$O(n)$ \quad mem $n$}{%
Answers a subtree DP for \emph{every} root at once, in two passes: up to collect each
subtree, then down to hand each child what the rest of the tree contributes. Shown for
$\mathrm{ans}[u] = \sum_v \mathrm{dist}(u,v)$.\\[0.3ex]
The pattern generalises: whatever \code{dp} accumulates from children, in the second
pass give child \code{v} the value ``everything except \code{v}'s subtree'', which for
a sum is \code{dp[u] - contribution(v)} and for a max needs the top two children.}

\begin{lstlisting}
// ans[u] = sum over v of dist(u, v)
vi sumOfDists(const vector<vector<int>>& g) {
    int n = (int)g.size();
    vi sub(n, 1), dp(n, 0), ans(n, 0);
    vector<int> par(n, -1), it(n, 0), st{0}, ord;
    while (!st.empty()) {
        int u = st.back();
        if (it[u] < (int)g[u].size()) {
            int v = g[u][it[u]++];
            if (v == par[u]) continue;
            par[v] = u; st.push_back(v);
        } else { ord.push_back(u); st.pop_back(); }
    }
    for (int u : ord) if (par[u] >= 0) {     // up: children
        sub[par[u]] += sub[u];
        dp[par[u]]  += dp[u] + sub[u];
    }
    ans[ord.back()] = dp[ord.back()];        // root
    down(i, (ll)ord.size(), 0) {             // down: parents
        int u = ord[i];
        if (par[u] >= 0)
            ans[u] = ans[par[u]] + n - 2 * sub[u];
    }
    return ans;
}
\end{lstlisting}

\entry{Offline DFS + rollback}{$O((n + q) \cdot \mathrm{apply})$ \quad mem $n + q$}{%
Not a structure, a shape. Bucket every query at the vertex it belongs to, then run
\emph{one} DFS from the root holding a global structure that always describes the
current root-to-\code{u} path. Apply \code{u}'s updates on entry, answer \code{u}'s
queries, recurse, and on exit undo exactly what was applied --- so the structure is
correct at every vertex without ever being rebuilt. Needs invertible updates and
offline queries.\\[0.3ex]
\begin{ctab}{0.33}{0.62}
array / counter & subtract what you added\\
diff array by depth & \code{d[dep] += x}, \code{d[dep+len+1] -= x}, prefix-summed
along the path\\
\code{BIT} & \code{add(i, x)} then \code{add(i, -x)}\\
\code{RDSU} & \code{snap()} then \code{undo(t)}\\
\end{ctab}\\[0.3ex]
Replace the recursion over the tree with a recursion over a segment tree of
\emph{time} and the same shape becomes offline dynamic connectivity: put each edge on
the $O(\log q)$ nodes covering the interval it is alive, DFS that tree, \code{join} on
entry, \code{undo} on exit. That is the only place \code{RDSU} is really needed ---
with an array or a BIT the undo is just the negated update.}

%NOCOMPILE
\begin{lstlisting}
void dfs(int u, int par) {
    int t = snapshot();                // e.g. rdsu.snap()
    for (auto& upd : at[u]) apply(upd);          // enter
    for (auto& q : qs[u]) ans[q.id] = get(q);
    for (int v : g[u]) if (v != par) dfs(v, u);
    undo(t);                     // exit: exact inverse
}
\end{lstlisting}
